% !TeX program = pdflatex
% Continuous IEEE Transactions supplement assembled from the verified RQ1--RQ2
% and RQ3 standalone supplements without changing their experimental content.
\documentclass[journal]{IEEEtran}

\usepackage[T1]{fontenc}
\usepackage[utf8]{inputenc}
\usepackage{amsmath,amsfonts,amssymb}
\usepackage{array}
\usepackage{booktabs}
\usepackage{graphicx}
\usepackage{makecell}
\usepackage{multirow}
\usepackage{tabularx}
\usepackage{stfloats}
\usepackage{xcolor}
\usepackage{url}
\usepackage{placeins}

\newcolumntype{Y}{>{\centering\arraybackslash}X}
\graphicspath{{./}}
\hyphenation{block-chain crowd-sens-ing rep-u-ta-tion}

\begin{document}

\renewcommand{\thefigure}{S\arabic{figure}}
\renewcommand{\thetable}{S\Roman{table}}
\markboth{IRPP-TD Supplementary Experimental Evidence}{RQ1--RQ3 Supplement}

\section{Supplementary Experimental Evidence for RQ1 and RQ2}
\label{sec:supplement}

This supplement reports secondary evidence omitted from the main experiment
section for space. It uses the same frozen data, calibration split, parameters,
and 30 random seeds. RQ1 adds distributional and participation diagnostics;
RQ2 gives the complete cold-start comparison, participation replication,
exact-ABOD calibration, resource diagnostics, and right-censoring-aware
early-error recovery. No parameter is retuned on tasks 21--100 or on an RQ2
stress condition.

\subsection{Extended RQ1 Evidence}
\label{sec:supp-rq1}

\textit{Statistical units and normalization:}
The Climate, Traffic, and Water workloads each contain 100 ordered tasks at
target participation 27 and 39. Tasks 1--20 define coordinate-wise
$Q_{0.95}-Q_{0.05}$ scales and freeze the model; tasks 21--100 are the
evaluation units. IRPP-TD is rerun with 30 anchor-sampling seeds. Its interval
estimator resamples seeds and tasks hierarchically, whereas deterministic
baselines use a task bootstrap. Paired intervals resample identical task IDs
and preserve the direction IRPP-TD minus baseline.

\begin{figure*}[!b]
\centering
\includegraphics[width=\textwidth]{fig_rq1_diagnostics.pdf}
\caption{Supplementary RQ1 diagnostics. (a) Realized participants per task for
the three scenes at target 27/39 (C: Climate, T: Traffic, W: Water).
(b) Distribution of IRPP-TD NRMSE over 30 anchor-sampling seeds.
(c) Paired IRPP-TD-minus-CRH-N NRMSE with 95\% task-bootstrap intervals; values
below zero favor IRPP-TD. (d) Absolute IRPP-TD NRMSE across calibration,
early-test, and mature-test phases at target 27.}
\label{fig:supp-rq1-diagnostics}
\end{figure*}

\begin{table*}[!t]
\centering
\caption{Paired RQ1 Comparison with CRH-N at Both Participation Levels}
\label{tab:supp-rq1-paired}
\scriptsize
\setlength{\tabcolsep}{4.2pt}
\renewcommand{\arraystretch}{1.02}
\begin{tabular}{lccccc}
\toprule
Scene & Target $\bar n$ & IRPP-TD NRMSE & CRH-N NRMSE &
Paired $\Delta$NRMSE [95\% CI] & Relative reduction [95\% CI] \\
\midrule
Climate & 27 & 0.003143 & 0.006025 &
$-2.882\,[-3.492,-2.235]\times10^{-3}$ & 47.8\% [39.8,54.5] \\
Climate & 39 & 0.002485 & 0.004710 &
$-2.226\,[-2.744,-1.693]\times10^{-3}$ & 47.2\% [40.0,53.2] \\
Traffic & 27 & 0.000526 & 0.001147 &
$-6.208\,[-8.435,-4.031]\times10^{-4}$ & 54.1\% [42.8,62.6] \\
Traffic & 39 & 0.000295 & 0.000952 &
$-6.570\,[-8.135,-4.936]\times10^{-4}$ & 69.0\% [63.5,73.3] \\
Water & 27 & 0.001597 & 0.012244 &
$-1.065\,[-1.674,-0.333]\times10^{-2}$ & 87.0\% [67.3,91.4] \\
Water & 39 & 0.001265 & 0.018050 &
$-1.678\,[-2.834,-0.379]\times10^{-2}$ & 93.0\% [75.4,95.8] \\
\bottomrule
\end{tabular}
\end{table*}

Fig.~\ref{fig:supp-rq1-diagnostics}(a) confirms that target participation is
not treated as a fixed worker count: the realized task ranges are 10--56 at
target 27 and 11--67 at target 39. Panel (b) shows that anchor sampling does not
change the accuracy order. Across the six scene--participation workloads,
run-to-run NRMSE standard deviations range from $2.32\times10^{-6}$ to
$2.19\times10^{-4}$, and every held-out run returns a finite truth.

Table~\ref{tab:supp-rq1-paired} and Fig.~\ref{fig:supp-rq1-diagnostics}(c)
extend the main CRH-N comparison to target 39. Every paired interval remains
strictly below zero. Panel (d) complements the main task-level CDF by exposing
phase behavior at target 27. From early to mature testing, IRPP-TD NRMSE changes
from 0.00347 to 0.00293 on Climate, from 0.000691 to 0.000396 on Traffic, and
from 0.00140 to 0.00170 on Water. The Water value remains below its calibration
level of 0.00200, so the evidence supports overall maturation without claiming
monotone improvement in every scene.

\FloatBarrier
\subsection{Extended RQ2 Protocol and Worker-Level Evidence}
\label{sec:supp-rq2}

\textit{Balanced paired replay:}
For every workload and seed, worker IDs are sorted by stored SUMO participation,
split into ten adjacent activity strata, randomly permuted within each stratum,
and interleaved before assigning 50 stable high-quality (H) and 50 stable
low-quality (L) workers. The largest H--L difference in mean participation over
all paired runs is 0.48 tasks. Only report values are regenerated: task truths,
participating IDs, order, and dimension remain aligned with RQ1. H reports have
truncated $\pm1\%$ noise, whereas L reports corrupt every coordinate by an
independently signed $5\%$--$20\%$ deviation. Ordinary report realizations are
paired across policies and early-error conditions.

\begin{table}[!t]
\centering
\caption{Cold-Start Policy Contract}
\label{tab:supp-rq2-policies}
\scriptsize
\setlength{\tabcolsep}{2.3pt}
\renewcommand{\arraystretch}{1.05}
\begin{tabularx}{\columnwidth}{@{}lXX@{}}
\toprule
Policy & Angular reference & External reports while bootstrapping \\
\midrule
IRPP-HQ5 & Reputation-guided RABOD & At most five trusted reports \\
IRPP-HQ10 & Same as IRPP-HQ5 & At most ten; HQ5 is its prefix \\
ABOD-HQ5 & Same anchors; exact ABOD & Same five trusted reports \\
No-Extra & RABOD; temporary ordinary fallback & None \\
Random-5 & RABOD & Five reports from a half-H/half-L pool \\
\bottomrule
\end{tabularx}
\end{table}

\textit{Comparison and calibration contract:}
The five policies share tasks, ordinary reports, normalization, reputation
state, stopping rules, seeds, and tuning budget. ABOD-HQ5 changes only the
angular evaluator, although its cutoffs must be calibrated separately because
exact-ABOD and sampled-RABOD scores have different scales. On all six RQ1
workloads and tasks 1--20 only, recomputation exactly reproduces the frozen
RABOD cutoffs $\mu_2=0.00172272$ and $\mu_1=0.0582259$ from 3,859 scores. The
same $3\times3$ quantile grid at frozen $\theta=0.35$ gives exact-ABOD cutoffs
$0.0217268$ and $0.0627577$. No RQ2 outcome is used for calibration.

\textit{Worker decisions and recovery:}
After at least $e_b=5$ post-prior outcomes, a worker is classified H when
$\psi_i\geq\theta$ and $\alpha_i^H-\alpha_i^L\geq q_b=2$, and L when the
reverse evidence gap is at least $q_b$; otherwise the system abstains.
Worker Macro-F1 retains abstentions as missed class decisions. E1 inserts one
false-low update at an H worker's second own participation; E3 inserts three at
participations 2--4. Stable recovery is the first completion of three
consecutive H decisions after the final forced error. Unresolved events at task
100 are right-censored.

\begin{figure*}[!t]
\centering
\includegraphics[width=\textwidth]{fig_rq2_coldstart_supplement_2x2.pdf}
\caption{Secondary RQ2 evidence. (a) Ten-task-block NRMSE for all five
cold-start policies at $\bar n=27$. (b) End-to-end NRMSE replication at target
participation 27 and 39. (c) Worker Macro-F1 under E0/E1/E3. (d)
Right-censoring-aware stable-recovery curves complementary to panel (c).
Curves and bands are scene-macro means and 95\% seed-bootstrap intervals over
30 paired seeds; lower is better in (a,b), whereas higher is better in (c,d).}
\label{fig:supp-rq2-worker}
\end{figure*}

\begin{table*}[!t]
\centering
\caption{Complete Cold-Start Outcomes and Participation Replication}
\label{tab:supp-rq2-cold}
\scriptsize
\setlength{\tabcolsep}{3.0pt}
\renewcommand{\arraystretch}{1.03}
\begin{tabular}{lcccccc}
\toprule
Policy & NRMSE-27 [95\% CI] & NRMSE-39 & Final F1 / cov. &
Pairs/task & Extra/task & Total extra \\
\midrule
IRPP-HQ5 & 0.001603 [0.001589,0.001616] & 0.001328 &
0.99994 / 0.99989 & 430 & 0.904 & 90.4 \\
IRPP-HQ10 & 0.001568 [0.001554,0.001581] & 0.001317 &
1.00000 / 1.00000 & 465 & 1.783 & 178.3 \\
ABOD-HQ5 & 0.001704 [0.001690,0.001720] & 0.001384 &
0.99788 / 0.99589 & 2,292 & 1.008 & 100.8 \\
No-Extra & 0.02534 [0.02434,0.02614] & 0.01798 &
0.358 / 0.925 & 411 & 0 & 0 \\
Random-5 & 0.02469 [0.02357,0.02565] & 0.01716 &
0.367 / 0.927 & 426 & 0.792 & 79.2 \\
\bottomrule
\end{tabular}
\par\vspace{0.45ex}
\noindent\parbox{0.98\textwidth}{\scriptsize
NRMSE-27 is the full 100-task value; the main figure uses mature tasks 51--100
for the accuracy--cost panel. Values are scene-macro estimates and NRMSE-27
includes a 95\% seed-bootstrap interval. Extra/task and Total extra count
reports actually submitted, rather than the registered bootstrap pool.}
\end{table*}

Fig.~\ref{fig:supp-rq2-worker}(a) shows that the accuracy gap is present across
task blocks rather than in one endpoint. At block 30, IRPP-HQ5 and IRPP-HQ10
reach Macro-F1 0.898 and 0.904; ABOD-HQ5 reaches 0.851; and No-Extra and
Random-5 reach 0.269 and 0.254. The latter two finish near 0.36 despite coverage above 0.92, so
their failure is confident wrong separation rather than only abstention.
Table~\ref{tab:supp-rq2-cold} and panel (b) preserve the ordering at
$\bar n=39$: additional ordinary participation lowers NRMSE but does not replace
a trustworthy initial reference.

Relative to ABOD-HQ5, IRPP-HQ5 evaluates 81.2\% fewer angular pairs per task.
The paired ABOD-minus-IRPP mature-NRMSE difference is
$9.11[8.05,10.26]\times10^{-5}$, and its final Macro-F1 difference is
$-0.206[-0.304,-0.119]$ percentage points. Increasing the trusted cap from five
to ten improves mature NRMSE by only $6.22[2.97,10.72]\times10^{-6}$ while
adding 87.9 [86.2,89.5] external reports per run. These paired results support
HQ5 as the more economical operating point without treating the small HQ10
accuracy change as a large practical gain.

\begin{figure*}[!t]
\centering
\includegraphics[width=\textwidth]{fig_rq2_coldstart_resource_detail_1x2.pdf}
\caption{RQ2 resource diagnostics at $\bar n=27$. (a) Angular-pair
evaluations, measured kernel time, and total screening time normalized by
IRPP-HQ5. (b) Cumulative external reports; labels give the fraction of tasks on
which the bootstrap is active. Exact ABOD evaluates substantially more pairs,
whereas vectorization and RABOD sampling overhead make small-workload wall time
non-proportional to pair count.}
\label{fig:supp-rq2-resources}
\end{figure*}

Fig.~\ref{fig:supp-rq2-resources} separates the implementation-independent
angular work from the present Python timing. Exact ABOD evaluates 5.33 times as
many pairs as IRPP-HQ5, but its vectorized kernel and total screening times are
1.01 and 0.65 times the IRPP-HQ5 values at $\bar n=27$. Therefore, RQ2 supports
lower sampled angular work, not a universal small-$n$ wall-clock speedup.
Runtime scaling and the crossover regime are evaluated in RQ4.

Fig.~\ref{fig:supp-rq2-worker}(c,d) distinguishes population classification
from event completion. After 10 own participations, 88.3\% of E1 workers and
53.5\% of E3 workers have recovered; after 20, the fractions are 99.3\% and
92.9\%. E3's paired NRMSE penalty peaks at $2.93\times10^{-4}$ in tasks 11--20
and then declines. Thus, one early mistake is usually absorbed quickly, while
three mistakes produce a measurable tail and 4.96\% right-censoring; the result
does not support instantaneous or universal recovery.

\textit{Paired early-error outcomes:}
Relative to E0, the full-horizon $\Delta$NRMSE values for E1 and E3 are
$0.061[0.010,0.111]\times10^{-4}$ and
$0.854[0.724,0.984]\times10^{-4}$; final Macro-F1 changes by $-0.04$ and
$-1.07$ percentage points. The restricted mean recovery delays are
7.26 [7.13,7.38] and 11.26 [11.11,11.41] own participations, with 0.24\% and
4.96\% right-censoring; the corresponding uncensored medians are 6 and 10.

\textit{Ablation semantics:}
The main-table Full row combines ternary TV-uncertainty reputation,
reputation-guided anchors, and adaptive seeds. No-U removes uncertainty
shrinkage; Binary-Rep. replaces the ternary state by a binary Beta state;
All-Anchors uses all current reports as anchors; Sequential performs all-worker
angular detection before reputation gating; No-Seeds deletes the bootstrap;
and Always-Seeds retains five external reports on every task. These changes
isolate reputation representation, reference construction, operation order,
bootstrap necessity, and adaptive switch-off cost, respectively.

\textit{Metric semantics:}
Screening-F1 compares per-report low decisions with the corruption mask recorded
before a method runs. Worker Macro-F1 instead evaluates the persistent H/L
state for all 100 IDs and penalizes abstention through recall. Clean retention
is the fraction of clean ordinary reports kept for TD, while bad leakage is the
fraction of bad ordinary reports retained. Anchor purity excludes external
workers. NRMSE normalizes each physical coordinate by its calibration-only
robust scale before cross-scene averaging. Extra/task counts reports actually
submitted, not registered external IDs.

\subsection{Reproducibility and Scope}
\label{sec:supp-reproducibility}

The formal run contains 1,080 run-level combinations, 108,000 task records,
and 99,000 worker-event records. It uses 30 paired seeds and 2,000
seed-bootstrap repetitions. Every IRPP-HQ5/HQ10 task returns a finite truth.
ABOD-HQ5 filters all reports on two Water tasks among its 18,000 evaluated
tasks (0.011\%); these no-truth outcomes remain in the reported rate rather than
being imputed. No task lacks angular input and no valid error is non-finite. The
recorded environment is Python 3.9.13 with NumPy 2.0.2, pandas 2.3.3, and SciPy
1.13.1. Configuration, source, six input JSON files, and formal results are
linked through SHA-256 hashes in the artifact manifest.

The conclusions are conditional on the preserved SUMO participation histories,
three task-aligned vector workloads, stated report-error process, disjoint
bootstrap workers, and a 100-task horizon. It does not establish robustness to
collusion, reputation building, seed corruption, scalar angular inputs, or a
malicious DR. Those properties remain assigned to RQ3--RQ5 and to the
protocol-level threat model.

\FloatBarrier
\section{Supplementary Experimental Evidence for RQ3}

\let\originalsection\section
\renewcommand{\section}{\subsection}
\section{Evaluation Contract and Trust Boundary}
\label{app:rq3-contract}

This supplement records the complete RQ3 contract behind Fig.~3 in the main
experimental section. RQ3 tests a worker-side adversary that can coordinate
ordinary worker reports and time its behavior using public task progress. The
adversary does not control the data requester, reputation blockchain, or trusted
bootstrap source. The $s_0=5$ bootstrap workers (IDs 1001--1005) are
platform-affiliated workers with aligned incentives and are assumed fully
honest. Their corruption is outside the worker-side adversarial model considered
in RQ3; no result below should be read as robustness to a compromised bootstrap
source.

The test reuses the RQ1 Climate, Traffic, and Water task incidence generated by
the stored SUMO trajectories. The primary workload has target participation
$\bar n=27$, and $\bar n=39$ is a replication. Worker IDs remain 1--100 and
each workload contains 100 tasks. Tasks 1--20 retain their calibration role;
accuracy is evaluated on tasks 21--100, except that attack-mode outcomes are
restricted to the blocks in which the corresponding attack is active. Table~
\ref{tab:rq3-settings} gives the complete fixed configuration.

\begin{table*}[!t]
\centering
\caption{Complete RQ3 Configuration}
\label{tab:rq3-settings}
\scriptsize
\setlength{\tabcolsep}{4.2pt}
\renewcommand{\arraystretch}{1.09}
\begin{tabular}{p{0.18\textwidth}p{0.76\textwidth}}
\toprule
Item & Fixed setting \\
\midrule
Workloads & Climate, Traffic, and Water; 100 tasks; target $\bar n=27$ in the main test and $\bar n=39$ for replication; original SUMO incidence retained. \\
Ordinary workers & IDs 1--100; malicious sets are activity-balanced, nested within each workload and seed, and contain $\operatorname{round}(100\rho_m)$ IDs. \\
Trusted bootstrap & IDs 1001--1005; honest, disjoint from ordinary workers, and supplied adaptively only while fewer than $n_b=5$ mature ordinary anchors exist. \\
Randomness & 30 paired seeds, 20260808--20260837; the same task incidence, honest report, malicious membership, attack report, and attack direction are replayed across compared methods. \\
Prevalence and strength & $\rho_m\in\{0,0.1,\ldots,0.8\}$ and $\kappa\in\{0.1,0.25,0.5,1,2\}$; the primary attack point is $(\rho_m,\kappa)=(0.3,0.5)$. \\
Schedules & Immediate attacks begin at task 1; On--Off attacks are active on 21--30, 41--50, 61--70, and 81--90; Mature-Anchor attackers behave honestly through task 40 and poison from task 41 onward. \\
Frozen IRPP-TD & $\zeta=0.25$, $\theta=0.35$, $\mu_1=0.05822585$, $\mu_2=0.00172272$, $e_b=n_b=s_0=5$, $q_b=2$, $\delta_{\max}=20$, $\epsilon_d=\epsilon_w=10^{-12}$, TD $\epsilon=10^{-5}$, and $t_{\max}=50$. \\
Screening thresholds & IRPP-TD uses its frozen RABOD decision and RPPS-TDC keeps the retained source threshold $p=0.3$. PRTD has no hard-rejection output, so its effective-weight gate is selected on a separate clean replay, tasks 1--20, to 96\% clean acceptance and then frozen. No attack outcome tunes any threshold. \\
Statistics & Scene-macro point estimates; 2,000 paired-seed bootstrap repetitions for 95\% intervals. No RQ3 outcome is used to retune a parameter. \\
Operational failure & $R_E\geq5$ or no-truth rate above 5\%, where $R_E$ is attacked NRMSE divided by the phase-matched clean NRMSE. This is an empirical operating criterion, not a theoretical breakdown point. \\
\bottomrule
\end{tabular}
\end{table*}

\section{Controlled Report and Attack Construction}
\label{app:rq3-attacks}

Let $\mathbf y_t$ be the task reference and let $\mathbf c$ and $\mathbf s$
be the coordinate-wise center and $Q_{0.95}-Q_{0.05}$ scale fitted only to the
RQ1 calibration partition. An honest coordinate is drawn from a zero-mean
Gaussian with standard deviation $0.01|y_{t,j}|/3$ and clipped between
$\min(0.99y_{t,j},1.01y_{t,j})$ and
$\max(0.99y_{t,j},1.01y_{t,j})$. We write
$\mathbf z=(\mathbf d-\mathbf c)/\mathbf s$ for the normalized report. Every
active attack satisfies, up to the stated compact jitter,
\begin{equation}
 \frac{1}{\sqrt{\ell}}
 \left\|\frac{\mathbf d^{\rm atk}-\mathbf d^{\rm hon}}{\mathbf s}\right\|_2
 =\kappa .
 \label{eq:rq3-strength}
\end{equation}

\textit{Independent bias:}
Each active malicious worker draws a deterministic seed-paired Rademacher
direction $\mathbf u_{i,t}\in\{-1,+1\}^{\ell}$ and submits
$\mathbf z^{\rm atk}_{i,t}=\mathbf z^{\rm hon}_{i,t}+\kappa\mathbf u_{i,t}$.

\textit{Compact collusion:}
The coalition computes the coordinate-wise median $\widetilde{\mathbf z}_t$ of
its own benign observations and submits reports around
$\widetilde{\mathbf z}_t+\kappa\mathbf v_1$, where $\mathbf v_1$ is the first
principal direction from calibration reports and
$\|\mathbf v_1\|_2=\sqrt{\ell}$. Independent Gaussian jitter has standard
deviation $0.02\kappa$. Thus, the attack is compact and coordinated but its
target does not read the task reference.

\textit{On--Off and Mature-Anchor poisoning:}
Both use the compact report when active. On--Off alternates the four active
blocks in Table~\ref{tab:rq3-settings} with benign blocks. Mature-Anchor behaves
honestly for 40 tasks so that coalition members can enter the reputation-guided
reference, then switches permanently at task 41. The latter directly probes the
reviewer-identified feedback loop rather than corrupting the trusted HQ5.

\section{Compared Implementations and Metrics}
\label{app:rq3-fairness}

IRPP-TD, CRH-N, PRTD, QE, and RPPS-TDC receive byte-identical ordinary reports
for a given replay. Panels (a) and (b) use the unchanged RQ1 aggregation
adapters and parameters. Fig.~3(c) compares method-faithful hard decisions:
IRPP-TD uses its native RABOD label, RPPS-TDC uses its native sequential
quality rule at $p=0.3$, and PRTD screens by its final effective weight
$n\gamma_i\omega_i/\sum_j\gamma_j\omega_j$ at the clean fourth-percentile
gate, then reruns the unchanged aggregation on retained reports. Clean
calibration is repeated inside each paired workload--seed replay, so its
variation enters the reported intervals. CRH-N and PRTD otherwise operate in
their native representation, and no task reference enters an aggregation
update. IRPP-TD additionally invokes its declared adaptive HQ5 bootstrap
mechanism. This is a system-level comparison rather than an equal-report-cost
comparison; RQ2 separately quantifies the bootstrap report cost.

The primary metric is returned-task NRMSE. No-truth is reported separately and
is never silently removed from the denominator of its rate. Malicious-report
leakage is the fraction of active malicious reports retained by IRPP-TD's
native RABOD decision, RPPS-TDC's native $p=0.3$ decision, or PRTD's
clean-calibrated relative-weight gate. Honest-report acceptance reports the
corresponding collateral cost. CRH-N has no report-selection output; QE's
exhaustive local-support scan remains in the accuracy comparison but is not
treated as the reputation-threshold filter in Fig.~3(c). Honest false-low is
the fraction of noncoalition reports assigned
the low-quality label. Ordinary-anchor purity excludes HQ5 and treats a
coalition member as malicious even during its reputation-building phase. This
definition makes the latent feedback risk visible before the task-41 switch.

\section{Complete RQ3 Results}
\label{app:rq3-results}

The following figure and tables report the replication, operating-boundary,
and clean-calibrated baseline evidence underlying the main-text panels.

\begin{figure*}[!t]
\centering
\includegraphics[width=0.95\textwidth]{fig_rq3_supplement_2x2.pdf}
\caption{Worker-level and replication evidence complementary to main-text
Fig.~3. (a) Compact-collusion prevalence at $\bar n=39$. (b) Mature-Anchor
error ratio over the complete prevalence--strength grid; crosses mark the
operational criterion. (c) IRPP-TD honest false-low rate across attack modes at
$(\rho_m,\kappa)=(0.3,0.5)$. (d) Full and No-U ordinary-anchor purity under
task-41 Mature-Anchor poisoning. Bars and bands carry 95\% paired-seed
bootstrap intervals.}
\label{fig:rq3-supplement}
\end{figure*}

\begin{table*}[!t]
\centering
\caption{Active-Task NRMSE Across Clean and Attack-Mode Replays at
$\bar n=27$}
\label{tab:rq3-mode-methods}
\scriptsize
\setlength{\tabcolsep}{8.0pt}
\renewcommand{\arraystretch}{1.05}
\begin{tabular}{lrrrrr}
\toprule
Method & Clean & Independent & Compact & On--Off & Mature-Anchor \\
\midrule
IRPP-TD & 0.00113 & 0.00137 & 0.00137 & 0.03631 & 0.03315 \\
CRH-N   & 0.00341 & 0.09963 & 0.04045 & 0.03544 & 0.04094 \\
PRTD    & 0.00127 & 0.00886 & 0.04634 & 0.04213 & 0.04430 \\
QE      & 0.00107 & 0.05216 & 0.14769 & 0.14883 & 0.14866 \\
\bottomrule
\end{tabular}
\par\vspace{0.4ex}
\parbox{0.91\textwidth}{\scriptsize Attack replays use
$(\rho_m,\kappa)=(0.3,0.5)$; Clean uses $\rho_m=0$. Each attack is evaluated
only on its active task set. Fig.~3(b) reports the corresponding 95\%
paired-seed bootstrap intervals.}
\end{table*}

\begin{table*}[!t]
\centering
\caption{Primary and Boundary RQ3 Outcomes}
\label{tab:rq3-complete}
\scriptsize
\textbf{(a) Compact-collusion prevalence: scene-macro NRMSE}\par\vspace{0.35ex}
\setlength{\tabcolsep}{4.0pt}
\begin{tabular}{ccrrrr}
\toprule
$\bar n$ & $\rho_m$ & IRPP-TD & CRH-N & PRTD & QE \\
\midrule
27 & 0.0 & 0.00113$\pm$0.00001 & 0.00341$\pm$0.00002 & 0.00127$\pm$0.00001 & 0.00107$\pm$0.00001 \\
27 & 0.3 & 0.00137$\pm$0.00002 & 0.0404$\pm$0.0036 & 0.0463$\pm$0.0038 & 0.148$\pm$0.001 \\
27 & 0.5 & 0.00164$\pm$0.00002 & 0.270$\pm$0.005 & 0.313$\pm$0.003 & 0.254$\pm$0.000 \\
27 & 0.8 & 0.00215$\pm$0.00002 & 0.465$\pm$0.001 & 0.488$\pm$0.000 & 0.406$\pm$0.000 \\
\midrule
39 & 0.0 & 0.00093$\pm$0.00001 & 0.00305$\pm$0.00003 & 0.00102$\pm$0.00001 & 0.00088$\pm$0.00001 \\
39 & 0.3 & 0.00112$\pm$0.00001 & 0.0196$\pm$0.0050 & 0.0305$\pm$0.0022 & 0.147$\pm$0.000 \\
39 & 0.5 & 0.00133$\pm$0.00001 & 0.259$\pm$0.004 & 0.306$\pm$0.003 & 0.252$\pm$0.001 \\
39 & 0.8 & 0.00201$\pm$0.00003 & 0.464$\pm$0.001 & 0.488$\pm$0.000 & 0.404$\pm$0.000 \\
\bottomrule
\end{tabular}
\par\vspace{1.1ex}
\begin{minipage}[t]{0.31\textwidth}
\centering\textbf{(b) Mature-anchor operational boundary}\par\vspace{0.35ex}
\setlength{\tabcolsep}{3.2pt}\begin{tabular}{rrrr}
\toprule $\kappa$ & Fail $\rho_m$ & $R_E@.8$ & Max NT (\%) \\
\midrule 0.1&0.3&141.05&0.02\\ 0.25&0.2&354.65&0.15\\ 0.5&0.2&710.72&0.22\\ 1&0.2&1422.11&0.13\\ 2&0.1&2844.40&0.19\\
\bottomrule\end{tabular}
\end{minipage}\hfill
\begin{minipage}[t]{0.31\textwidth}
\centering\textbf{(c) Malicious-report leakage}\par\vspace{0.35ex}
\setlength{\tabcolsep}{3.0pt}\begin{tabular}{lrrr}
\toprule Mode & IRPP-TD & RPPS-TDC & PRTD \\
\midrule Independent&0.000&0.00025&0.719\\ Compact&0.000&0.080&0.777\\ On--Off&0.249&0.293&0.780\\ Mature&0.206&0.295&0.779\\
\bottomrule\end{tabular}
\end{minipage}\hfill
\begin{minipage}[t]{0.31\textwidth}
\centering\textbf{(d) Ordinary-anchor purity}\par\vspace{0.35ex}
\setlength{\tabcolsep}{3.2pt}\begin{tabular}{lrrr}
\toprule Variant & B40 & B50 & B100 \\
\midrule Full&0.700&0.710&0.968\\ Binary-Rep.&0.700&0.711&0.984\\ No-U&0.700&0.712&0.986\\ All-Anchors&0.699&0.699&0.697\\ Sequential&0.699&0.699&0.697\\
\bottomrule\end{tabular}
\end{minipage}
\par\vspace{0.4ex}
\parbox{0.97\textwidth}{\scriptsize Panel (a) reports means $\pm$ the larger
95\% paired-seed bootstrap half-width. Failure in (b) means $R_E\geq5$ or
no-truth above 5\%. Thresholds in (c) are fixed before attack evaluation;
only PRTD's added gate is fitted on clean tasks 1--20.}
\end{table*}

Table~\ref{tab:rq3-complete}(a) and Fig.~\ref{fig:rq3-supplement}(a)
show that the immediate compact-collusion conclusion replicates at the larger
participation level. At $\bar n=27$ and
$\rho_m=0.8$, IRPP-TD has 0.00215 NRMSE, whereas CRH-N, PRTD, and QE have
0.465, 0.488, and 0.406. The prevalence sweep has zero no-truth outcomes for
all four methods.

Table~\ref{tab:rq3-mode-methods} complements the prevalence sweep with matched
algorithm comparisons across attack schedules. IRPP-TD is lowest under both
immediate modes, with 0.00137 NRMSE versus 0.00886 for the closest Independent
baseline and 0.04045 for the closest Compact baseline. QE is marginally lower
on clean data, at 0.00107 versus 0.00113. Under On--Off, CRH-N and IRPP-TD give
0.03544 and 0.03631; under Mature-Anchor, IRPP-TD is lowest at 0.03315.
Thus, reputation-aware timing narrows the immediate-attack advantage without
overturning IRPP-TD's unique near-clean result across both immediate geometries.

The delayed grid in Table~\ref{tab:rq3-complete}(b) exposes a different operating
regime. Mature attackers cross the $R_E=5$ criterion at smaller prevalence as
$\kappa$ increases. No-truth remains below 0.23\% in every cell; hence these
failures reflect inaccurate returned aggregates, not widespread abstention.
The full surface is shown in Fig.~\ref{fig:rq3-supplement}(b). At the primary
point, IRPP-TD fully rejects immediate Independent and Compact attacks, whereas
On--Off and Mature-Anchor reports leak after first building reputation.
IRPP-TD leakage is 0, 0, 0.249, and 0.206 across Independent, Compact, On--Off,
and Mature-Anchor. Native RPPS-TDC gives 0.00025, 0.080, 0.293, and 0.295;
its honest-report acceptance is 0.730--0.986, below IRPP-TD's 0.953--0.993.
PRTD accepts over 0.999 of honest reports but also retains 0.719--0.780 of
poison. The RPPS-TDC source threshold accepts 52.9\% of reports on the separate
clean tasks 1--20, confirming that its low immediate leakage is accompanied by
aggressive collateral filtering. IRPP-TD therefore provides the strongest
joint accuracy, poison rejection, and honest-report preservation, while timed
attackers remain a documented feedback limitation.

Table~\ref{tab:rq3-complete}(d) and Fig.~\ref{fig:rq3-supplement}(d) separate
containment from prevention. Full,
Binary-Rep., and No-U guided anchoring all remove coalition members from the
ordinary anchor set after the switch; Full purity reaches 0.968 by block 100.
All-Anchors and Sequential cannot purify that reference and remain near the
population baseline of 0.70. For Full, the phase-matched error ratio falls from
71.61 in block 50 to 54.78, 28.64, 9.23, 3.39, and 1.35 in blocks
60--100. For On--Off, inactive blocks 40, 60, 80, and 100 return to ratios
1.03, 1.05, 1.06, and 1.07, but every later active block again incurs a large error.
Thus, reputation-guided anchoring supports recovery from a permanent switch;
without explicit forgetting or change-point logic, it does not prevent repeated
bursts from exploiting stale reputation.

\section{Reproducibility and Integrity Checks}
\label{app:rq3-reproducibility}

The isolated rerun does not modify the legacy IRPP source; six formal shards
each contain five seeds. The primary run records 1,044,000 task outcomes and
2,226,690 report decisions, while the mode extension and Mature-Anchor grid add
135,000 and 405,000 task outcomes. The clean-calibrated QE/PRTD audit adds
414,000 task outcomes and 180 workload--seed calibrations; the RPPS-TDC
extension adds 36,000 sequential screening records and 90 clean-threshold
measurements, and the no-gate diagnostic records 720 rates. Automated checks
verify all 30 seeds, unique keys, finite returned metrics, nested malicious
sets, honest HQ5, schedule transitions, PRTD clean acceptance within
0.95--0.97, fixed RPPS-TDC $p=0.3$, nondegenerate leakage, and active-poison
membership. All primary and extension manifests agree on generator, wrapper,
and adapter hashes. Retained-source SHA-256 prefixes are \texttt{95a89b80}
(CRH-N), \texttt{e4242705} (QE), \texttt{1f5a5d63} (PRTD),
\texttt{755e7d7c} (RPPS-TDC TD), and \texttt{6d3a27c9} (RPPS-TDC reputation);
the integrity-audit JSON records the complete check.
\let\section\originalsection

\end{document}
